\chapter{Conclusion}
\label{chap:conclusion}

\section{Summary of Contributions}
\label{sec:contributions}

This thesis developed an AI-driven framework for resource demand forecasting and allocation recommendation in CAPEX portfolio management. The principal contributions are:

\begin{enumerate}
  \item A temporal skill-demand forecasting model trained on historical CAPEX project data, capable of predicting skill-specific resource requirements across project lifecycle phases.
  \item A proactive bottleneck detection algorithm that identifies anticipated resource scarcities at sufficient lead time for corrective action.
  \item A constraint-aware optimization model that translates demand forecasts into actionable allocation recommendations aligned with organizational constraints.
  \item An integrated system architecture embedding the above components into a cohesive decision-support engine.
\end{enumerate}

% TODO: expand with quantitative highlights from evaluation

\section{Answers to Research Questions}
\label{sec:rq_answers}

% TODO: explicitly address each RQ from Section~\ref{sec:research_questions}

\section{Future Work}
\label{sec:future_work}

Several avenues remain open for future research:

\begin{itemize}
  \item \textbf{Real-time adaptation:} Extending the framework to incorporate live project status updates and dynamically re-optimize allocations.
  \item \textbf{Multi-objective optimization:} Incorporating additional objectives such as cost minimization, employee workload balancing, or skill development goals.
  \item \textbf{Cross-portfolio learning:} Training models on data from multiple organizations to improve generalizability and reduce data requirements per site.
  \item \textbf{Natural language interfaces:} Enabling portfolio managers to query the system and receive recommendations through conversational AI interfaces.
\end{itemize}
